\documentclass[11pt]{article}

\usepackage{amsmath}
\usepackage{amssymb}
\usepackage{amsthm}
\usepackage{mathtools}
\usepackage{booktabs}
\usepackage[margin=1in]{geometry}
\usepackage[hidelinks]{hyperref}

\theoremstyle{definition}
\newtheorem{definition}{Definition}[section]

\theoremstyle{plain}
\newtheorem{theorem}{Theorem}[section]
\newtheorem{lemma}[theorem]{Lemma}
\newtheorem{proposition}[theorem]{Proposition}
\newtheorem{corollary}[theorem]{Corollary}

\theoremstyle{remark}
\newtheorem*{remark}{Remark}

\newcommand{\Star}[1]{{#1}^{*}}
\newcommand{\derives}{\Rightarrow}
\newcommand{\derivesstar}{\Rightarrow^{*}}
\newcommand{\Reg}{\textsf{Reg}}
\newcommand{\CF}{\textsf{CF}}
\newcommand{\CS}{\textsf{CS}}
\newcommand{\RE}{\textsf{RE}}

\title{Notes on Grammatical Analysis}
\author{Sinan Olsson-Pasic}
\date{January 9, 2026\\[2pt]\small(revised May 26, 2026)}

\begin{document}
\maketitle

The following are a collection of notes that I'm documenting during my
development and learning process of my grammatical analysis engine.

\section{What is a language?}

A language to a computer scientist is a possibly infinite set of sentences,
each made up of ``tokens'' (words), and each token made up of ``symbols''
(letters). Or in other words, a language is a subset $\Star{\Sigma}$ for some
alphabet $\Sigma$. Which ``sentences'' belong to a ``language'' and which
``tokens'' make up those ``sentences'' is explained by the definition of a
generative grammar for the language.

A generative grammar is a 4-tuple $(V_N, V_T, R, S)$ such that:
\begin{enumerate}
  \item $V_N$ and $V_T$ are finite sets of symbols.
  \item $V_N \cap V_T = \emptyset$
  \item $R$ is the set of pairs $(P, Q)$ such that:
    \begin{enumerate}
      \item $P \in \Star{(V_N \cup V_T)}{}^{+}$
      \item $Q \in \Star{(V_N \cup V_T)}$
    \end{enumerate}
  \item $S \in V_N$
\end{enumerate}

Where $V_N$ is the set of non-terminals, $V_T$ is the set of terminals, $R$ is
the set of rules, and $S$ is the start symbol. $\emptyset$ denotes the empty
set, meaning the intersection of the terminals and nonterminals returns no
element.

A language $L$ bound by a grammar $G$ therefore is defined as:
\begin{equation}
  L(G) = \{\, w \in \Star{V_t} \mid S \derivesstar w \,\}
\end{equation}

All of this is to say that a generative grammar gives us the specification of
how sentences for a language is constructed rather than listing the valid
sentences. This allows us to describe infinite languages using finite
information.

\subsection{On the derivation relation}

The arrow $\derivesstar$ hides the mechanism of generation, so it's worth
unfolding. A rule $(P, Q) \in R$ licenses one rewriting step: a string
$\alpha P \beta$, with the pattern $P$ sitting in some context $\alpha, \beta$,
becomes $\alpha Q \beta$. Call this one-step relation $\derives$.

\begin{definition}[Derivation]
$\derives \, \subseteq \Star{(V_N \cup V_T)} \times \Star{(V_N \cup V_T)}$ is
the relation with $\alpha P \beta \derives \alpha Q \beta$ whenever
$(P,Q) \in R$. Its reflexive--transitive closure is $\derivesstar$:
$\alpha \derivesstar \gamma$ holds iff $\alpha = \gamma$ or there is a finite
chain $\alpha = \gamma_0 \derives \gamma_1 \derives \cdots \derives \gamma_k =
\gamma$.
\end{definition}

So $\derivesstar$ is $\derives$ applied any finite number of times, and $L(G)$
collects exactly the terminal strings reachable from $S$ this way. Membership of
$w$ in $L(G)$ is an existential statement: it asserts that \emph{some}
derivation exists. The same $w$ may have many derivations, or none. A
derivation is a proof of membership; its absence is the stronger claim that no
proof exists.

\subsection{A note on infinite languages and descriptions}

Is it possible to give a description for every language? The answer is no. The
reason is a counting argument.

\begin{theorem}[Almost every language is indescribable]
Over a finite alphabet $\Sigma$ there are uncountably many languages but only
countably many grammars. So the languages that admit any finite description are
a vanishing fraction of all languages.
\end{theorem}

\begin{proof}[Proof sketch]
$\Star{\Sigma}$ is countable: list strings by length, then alphabetically, for a
bijection with $\mathbb{N}$, so $|\Star{\Sigma}| = \aleph_0$. A language is any
subset, so the languages form the power set, and
$|\mathcal{P}(\Star{\Sigma})| = 2^{\aleph_0} > \aleph_0$ by Cantor's theorem ---
uncountably many. (Diagonally: any list $L_0, L_1, \dots$ misses
$D = \{ w_i \mid w_i \notin L_i \}$, which disagrees with each $L_i$ on the
$i$-th string $w_i$.) A grammar is a finite object --- finite $V_N, V_T, R$ ---
hence writable as a finite string, and the finite strings over a fixed alphabet
are countable. Countably many descriptions cannot name uncountably many
languages.
\end{proof}

The same bound applies to Turing machines, regular expressions, or descriptions
of any finitary kind. So the grammar hierarchy below sits entirely inside the
countable sliver of describable languages. The question is never whether an
arbitrary language can be described, but how much descriptive power natural
languages demand, and what that power costs.

\section{The Chomsky hierarchy}

The 4-tuple above places no restriction on the shape of a rule: any nonempty
$P$ may rewrite to any $Q$. Restricting that shape carves out a sequence of
grammar classes, each weaker than the last and each cheaper to decide.

\begin{definition}[The four grammar types]
A grammar $(V_N, V_T, R, S)$ is:
\begin{enumerate}
  \item \textbf{Type 0 (recursively enumerable).} No restriction: any $(P,Q)$
        with $P \in \Star{(V_N \cup V_T)}{}^{+}$.
  \item \textbf{Type 1 (context-sensitive).} Every rule is non-contracting,
        $|P| \le |Q|$; equivalently of the form
        $\alpha A \beta \to \alpha \gamma \beta$ with $A \in V_N$,
        $\gamma \neq \varepsilon$.
  \item \textbf{Type 2 (context-free).} Every left-hand side is a single
        nonterminal: $A \to Q$.
  \item \textbf{Type 3 (regular).} Every rule is right-linear: $A \to aB$ or
        $A \to a$.
\end{enumerate}
\end{definition}

The names are literal. Context-free: a nonterminal rewrites regardless of its
neighbors. Context-sensitive: the frame $\alpha\,\underline{\phantom{A}}\,\beta$
must be present for the rewrite to fire. The content is not the rule syntax but
the machine each class corresponds to.

\begin{table}[h]
\centering
\begin{tabular}{@{}llll@{}}
\toprule
Type & Class & Recognizer & Memory \\
\midrule
0 & recursively enumerable & Turing machine & unbounded tape \\
1 & context-sensitive & linear-bounded automaton & tape linear in input \\
2 & context-free & pushdown automaton & a single stack \\
3 & regular & finite automaton & none \\
\bottomrule
\end{tabular}
\caption{Grammar power tracks memory.}
\end{table}

Each grammar class is generated exactly by machines with a fixed memory budget.
No memory tracks only a bounded amount of information; one stack matches nested
brackets but cannot compare three counts at once; unbounded tape recognizes
anything recognizable. The classes nest strictly:
\[
  \Reg \subsetneq \CF \subsetneq \CS \subsetneq \RE,
\]
with $\{a^n b^n\} \in \CF \setminus \Reg$ and
$\{a^n b^n c^n\} \in \CS \setminus \CF$. The containments are immediate from the
definitions; the strictness is what the pumping lemmas of Section~4 establish.
These two sets are the abstract shapes that natural-language dependencies turn
out to take.

\section{The membership problem}

The question a grammar checker actually asks is membership: given $w$, is
$w \in L(G)$? The hierarchy answers it at every rung, and the answers set the
budget for any checker.

\begin{theorem}[Cost of membership]\label{thm:cost}
For fixed $G$ and input length $n$: regular is $O(n)$; context-free is $O(n^3)$;
context-sensitive is \textsf{PSPACE}-complete; type 0 is undecidable.
\end{theorem}

\begin{proof}[Proof sketches]
\emph{Regular.} A DFA reads each symbol once; acceptance is a final table
lookup.

\emph{Context-free.} CKY (in Chomsky normal form, $A \to BC$ or $A \to a$) fills
a triangular table $T[i,j]$ of all nonterminals deriving the length-$j$
substring at position $i$: $O(n^2)$ cells, each tried against $O(n)$ split
points, so $O(n^3 |G|)$. Then $w \in L(G)$ iff $S \in T[1,n]$. The table also
records every parse, not just the verdict. (Valiant pushes the constant toward
matrix-multiplication time, $\approx O(n^{2.37})$.)

\emph{Context-sensitive.} A linear-bounded automaton uses tape linear in $n$, so
each configuration fits in polynomial space; configuration reachability is in
\textsf{PSPACE}, with a matching reduction for completeness.

\emph{Type 0.} A derivation can simulate an arbitrary Turing machine, so
deciding $S \derivesstar w$ would decide halting.
\end{proof}

A fully general grammar is therefore not merely expensive but unanswerable. A
workable checker has to sit at or near context-free: expressive enough for real
structure, cheap enough to decide in polynomial time. One distinction is needed
before going further.

\begin{definition}[Weak vs.\ strong generative capacity]
The \emph{weak} capacity of $G$ is the string set $L(G)$; the \emph{strong}
capacity is the set of parse trees it assigns. Grammars are weakly equivalent
when they generate the same strings, even with different trees.
\end{definition}

Deciding ``grammatical?'' is a weak-capacity question. Locating and naming an
error is a strong-capacity one: pointing at ``the verb phrase'' needs the right
tree, not just the right yes/no.

\section{Where natural language sits}

Two results fix the position of natural language. Both are pumping arguments,
so state the tools first.

\subsection{Not regular}

\begin{lemma}[Regular pumping]
If $L$ is regular, some $p \ge 1$ has: every $s \in L$ with $|s| \ge p$ factors
as $s = xyz$, $|xy| \le p$, $|y| \ge 1$, and $x y^i z \in L$ for all $i \ge 0$.
\end{lemma}

\begin{proof}[Proof sketch]
Let $p$ be the number of DFA states. A string of length $\ge p$ revisits some
state (pigeonhole); the loop $y$ between the two visits can be traversed any
number of times while staying on an accepting run.
\end{proof}

\begin{proposition}
$\{a^n b^n\}$ is not regular.
\end{proposition}

\begin{proof}
With $s = a^p b^p$ and $|xy| \le p$, the block $y$ is all $a$'s, so $xy^2z$ has
more $a$'s than $b$'s and leaves the language.
\end{proof}

Center embedding gives this shape in natural language: ``the rat [the cat [the
dog chased] killed] ate'' nests matching dependencies that count like
$a^n b^n$. A finite-state device has no stack to count nesting depth, so natural
language needs at least a stack --- at least context-free (Chomsky, 1956).

\subsection{Not context-free}

\begin{lemma}[Context-free pumping]
If $L$ is context-free, some $p \ge 1$ has: every $z \in L$ with $|z| \ge p$
factors as $z = uvwxy$, $|vwx| \le p$, $|vx| \ge 1$, and $u v^i w x^i y \in L$
for all $i \ge 0$.
\end{lemma}

\begin{proof}[Proof sketch]
A long enough $z$ forces a root-to-leaf path in its parse tree longer than
$|V_N|$, so some nonterminal repeats on that path. The subtree between the two
occurrences can be excised or duplicated, pumping the frontier pieces $v$ and
$x$ together; capping the lower occurrence's height bounds $|vwx|$.
\end{proof}

\begin{proposition}
$\{a^n b^n c^n\}$ and the copy language $\{ww\}$ are not context-free.
\end{proposition}

\begin{proof}
For $z = a^p b^p c^p$, the window $vwx$ spans at most two of the three blocks,
so pumping changes at most two counts and breaks the equality. The same
straddling argument breaks $\{ww\}$, which is the abstract form of a
\emph{crossing} rather than nesting dependency.
\end{proof}

Whether natural language stays context-free was open for decades and settled by
Shieber (1985), using cross-serial dependencies in Swiss German subordinate
clauses: a block of case-marked noun phrases followed by a block of verbs, where
each verb must match the correspondingly-placed noun phrase in case --- so the
dependencies cross.

\begin{theorem}[Shieber 1985]
Some natural language has a string set that is not context-free.
\end{theorem}

\begin{proof}[Proof architecture]
Context-free languages are closed under intersection with regular languages and
under homomorphism. For Swiss German $D$, a suitable regular $R$ and
homomorphism $h$ give
\[
  h(D \cap R) = \{\, a^m\, b^n\, c^m\, d^n \mid m,n \ge 0 \,\},
\]
where $a,b$ are the two case classes of noun phrase and $c,d$ the verbs that
govern them; crossing forces the interleaved matched counts. Pumping shows this
set is not context-free, and since $\CF$ is closed under the operations used,
$D$ cannot be context-free either.
\end{proof}

So natural language is strictly above context-free, but only just. Joshi (1985)
named the target class.

\begin{definition}[Mildly context-sensitive]
A language class is mildly context-sensitive if it contains the context-free
languages, captures limited cross-serial dependencies, has the constant-growth
property (the length set $\{|w| : w \in L\}$ is semilinear, ruling out
$\{a^{2^n}\}$), and is parsable in polynomial time.
\end{definition}

\begin{theorem}[Vijay-Shanker \& Weir 1994]
Tree-Adjoining Grammar, Combinatory Categorial Grammar, Linear Indexed Grammar,
and Head Grammar are weakly equivalent --- one class, properly above
context-free, inside the mildly context-sensitive region.
\end{theorem}

Four independently motivated formalisms generating one class is good evidence
the class is natural. The cost shows in parsing: TAG recognition is $O(n^6)$,
polynomial but visibly dearer than $O(n^3)$. In practice a context-free skeleton
plus features covers almost everything worth checking, and the genuinely
cross-serial constructions are rare enough to set aside.

\section{From generation to constraints}

The grammars above generate: start at $S$ and rewrite down to $w$. The
constraint-based frameworks (HPSG, LFG) invert this, a view called
model-theoretic syntax (Pullum \& Scholz, 2001): a grammar is not a device that
produces strings but a set of constraints a structure must satisfy.
Grammaticality becomes the existence of a structure for $w$ satisfying every
constraint.

A violated constraint is local --- it concerns particular constituents --- and
named --- it is a particular constraint. That is what makes error reporting
possible at all, rather than a bare yes/no. The machinery is feature structures
and unification.

\begin{definition}[Feature structure]
A rooted, connected, directed acyclic graph with feature-labelled arcs and
atomic leaf values, sharing of substructure allowed (reentrancy). Equivalently a
partial function from feature paths to values, e.g.
$[\, \textsf{cat}\!:\textsf{noun},\ \textsf{num}\!:\textsf{sg},\
\textsf{pers}\!:3,\ \textsf{case}\!:\textsf{nom}\,]$.
\end{definition}

\begin{definition}[Subsumption and unification]
$F \sqsubseteq G$ (\,$F$ subsumes $G$\,) when $G$ carries all of $F$'s
information and perhaps more --- $F$ is more general. The unification
$F \sqcup G$ is the least upper bound in this order: the most general structure
holding the information of both.
\end{definition}

\begin{theorem}[Unification is a lattice join]
Under subsumption, feature structures form a lattice once a top element $\top$
(inconsistency) is adjoined. $F \sqcup G$ always exists and equals $\top$
exactly when $F$ and $G$ conflict on some shared path.
\end{theorem}

\begin{proof}[Proof sketch]
Compatible information merges to a unique least refinement, computed by a graph
merge (union--find over the two DAGs, identifying nodes reached by the same
path). Two paths demanding distinct atomic values --- $\textsf{num}\!:
\textsf{sg}$ against $\textsf{num}\!:\textsf{pl}$ --- have no consistent
refinement, so the join is $\top$: unification fails.
\end{proof}

Agreement is enforced by unifying the relevant constituents' feature structures.
A subject with $\textsf{num}\!:\textsf{sg}$ against a verb with $\textsf{num}\!:
\textsf{pl}$ unifies to $\top$: the parse fails, and the failure says which two
constituents clashed and on which feature. That is the gap between a yes/no
acceptor and one that can report a verb disagreeing with its singular subject.
Kasper and Rounds (1986) gave this a logical semantics, so rules read as
constraints rather than ad hoc checks; the verdict and the diagnosis come from
one mechanism.

\section{Gradience and the coverage trap}

Two cracks in the purely formal program shaped the field afterward.

First, grammaticality is binary but acceptability is graded (the familiar ``?'',
``??'', ``*''). This drove weighted grammars. A probabilistic context-free
grammar attaches a probability to each production $A \to \beta$ with
$\sum_\beta P(A \to \beta) = 1$; a tree's probability is the product of its
productions, and probabilistic CKY returns the most likely tree. The field then
moved to data-driven statistical parsing and, now, neural sequence models;
modern grammatical error correction is sequence-to-sequence, not rule-based,
because it tolerates gradience and gaps.

Second, and crucially, the coverage caveat is formal, not incidental. A grammar
$G$ defines $L(G)$, but any hand-built $G$ is incomplete, so
\[
  \neg\big(w \in L(G)\big) \;\neq\; \text{``$w$ is ungrammatical.''}
\]
No parse may mean only that $G$ does not yet cover $w$. Treating ``no parse'' as
``error'' manufactures false positives. So a checker reads failure as
\emph{unrecognized}, not \emph{ungrammatical}: it reports violated constraints
it knows (with targeted error rules, or explicit mal-rules that match specific
wrong patterns), never the bare absence of a parse.

\section{Constraints on the rule design}

What the theory dictates, collected:
\begin{itemize}
  \item Anchor at the context-free level --- membership must be decidable and
        polynomial (Theorem~\ref{thm:cost}).
  \item Natural language exceeds context-free only mildly, on rare cross-serial
        constructions; defer them.
  \item Use feature structures and unification for agreement, government, and
        case; a unification failure already localizes and names the error.
  \item Scope the negative verdict honestly: report violated known constraints,
        never the absence of a parse.
\end{itemize}

\section{The engine as a pipeline}

Membership and diagnosis are computed in stages, each a function from one
representation to the next:

\begin{verbatim}
  text -> tokenizer -> morphology -> lexicon -> parser(+unification)
       -> error detection -> report
\end{verbatim}

\begin{enumerate}
  \item \textbf{Tokenizer.} Raw string $\to$ sequence of tokens.
  \item \textbf{Morphology.} Each token $\to$ (lemma, morphological features),
        possibly several analyses.
  \item \textbf{Lexicon.} Each analysis $\to$ a category and a feature
        structure.
  \item \textbf{Parser.} Tokens $\to$ constituent structure over a context-free
        backbone, unifying features as it builds (Section~3 fixes this at the
        context-free level for decidability).
  \item \textbf{Error detection.} Unification failures and matched error
        patterns $\to$ localized violations.
  \item \textbf{Report.} Violations $\to$ highlighted span, named rule,
        suggested fix.
\end{enumerate}

The pipeline is mostly language-independent; the lexicon, the morphology, and
the rule set are the language-specific data it runs on. Stages 1--3 are where
ambiguity enters --- a token may have several analyses --- and the parser
resolves it by keeping only the analyses that fit into a structure.

\section{Tokenization}

The tokenizer realizes a map $\Star{\Sigma} \to \Star{T}$ from the symbol string
to a token sequence over a vocabulary $T$. For space-delimited languages
(English, Swedish) whitespace and punctuation carry most of the work, but the
map is not trivial: contractions split (\texttt{don't} $\to$ \texttt{do} +
\texttt{n't}), clitics and hyphenates need a policy, and punctuation attaches
ambiguously. For languages written without spaces (Japanese) segmentation is the
hard part --- word boundaries are ambiguous and must be recovered by dictionary
longest-match or a statistical model. Tokenization can itself be ambiguous, so
it is better treated as a relation than a function, with the ambiguity passed
downstream rather than resolved prematurely.

\section{Morphological analysis}

A word form decomposes into morphemes: a stem plus affixes. Inflectional
morphology marks features --- number, person, tense, case, gender --- without
changing category; derivational morphology changes it. The analyzer maps a
surface form to one or more (lemma, features) pairs:

\begin{verbatim}
  bark  ->  (bark, V, pres, agr != 3sg)     # "they bark"
  bark  ->  (bark, N, num=sg)               # "a bark"
  barks ->  (bark, V, pres, agr  = 3sg)     # "it barks"
  barks ->  (bark, N, num=pl)               # "two barks"
\end{verbatim}

Morphology is regular (Type 3 in Section~2), so finite-state transducers handle
it --- Koskenniemi's two-level model, which is fast and runs in both directions,
analysis and generation. For English the morphology is light and the residual
ambiguity (is \texttt{bark} a noun or a verb?) is left for the parser. For
agglutinative (Japanese, Turkish) or fusional (Russian) languages this stage
carries most of the grammatical information, and the feature structures it emits
are correspondingly richer.

\section{Feature structures in the grammar}

Section~5 gives feature structures and unification abstractly; in the engine they
attach to two things. Lexicon entries carry one feature structure each. Grammar
rules carry \emph{equations} over the features of their constituents --- the
constraints. A small English fragment:

\begin{verbatim}
  S  -> NP VP    : NP.agr = VP.agr      # subject-verb agreement
  NP -> Det N    : NP.agr = N.agr
  VP -> V        : VP.agr = V.agr
  Det -> the
  N   -> dog     : [num=sg, pers=3]
  V   -> bark    : [agr != 3sg]
  V   -> barks   : [agr  = 3sg]
\end{verbatim}

A context-free backbone annotated with feature equations is a \emph{unification
grammar}; PATR-II is the minimal such formalism. As the parser applies a rule it
unifies the daughters' structures according to the rule's equations and derives
the mother's. A satisfied equation propagates information upward; a violated one
yields $\top$ and pins the failure to a specific rule and a specific pair of
constituents. The verdict and its explanation are then the same computation seen
two ways.

\section{Mal-rules and targeted error recognition}

The coverage trap (Section~6) forbids reading ``no parse'' as ``ungrammatical.''
Two strategies turn an absent parse into a specific, reportable diagnosis.

\textbf{Constraint relaxation.} When unification fails, retry the parse with one
constraint dropped. If relaxing exactly the agreement equation
\texttt{NP.agr = VP.agr} lets the parse complete, that equation is the
diagnosis and the span is the constituents it related. General, but it
over-triggers when several relaxations succeed at once.

\textbf{Mal-rules.} Encode a known wrong pattern as a rule carrying an error
label, an explanation, and a correction. From the language-learning literature,
mal-rules target specific recurrent errors:

\begin{verbatim}
  S -> NP VP  : NP.agr = 3sg, VP.agr != 3sg
              *ERR  subject-verb agreement (missing 3sg -s)
               FIX  inflect verb to 3sg, or pluralize subject
\end{verbatim}

Precise but anticipatory --- each error must be foreseen and written. A working
engine parses normally first and only on failure falls back to relaxation and
mal-rules, ranked by likelihood, reporting the \emph{minimal} set of violated
constraints rather than every possible complaint.

\section{A worked example: ``the dog bark''}

Trace the pipeline on the ungrammatical \texttt{the dog bark}.

\emph{Tokenize.} \texttt{[the, dog, bark]}.

\emph{Morphology and lexicon.} \texttt{the} is a determiner, unmarked for
number. \texttt{dog} is a noun, \texttt{[num=sg, pers=3]}. \texttt{bark} is
ambiguous: a present-tense verb requiring \texttt{agr != 3sg}, or a singular
noun.

\emph{Parse.} Reading \texttt{bark} as a verb:

\begin{verbatim}
        S
       / \
     NP   VP
    /  \   |
  Det   N  V
  the  dog bark
\end{verbatim}

The subject NP inherits the noun's features, \texttt{NP.agr = [num=sg, pers=3]}
(that is, 3sg). The VP inherits the verb's requirement, \texttt{VP.agr =
[agr != 3sg]}. The rule \texttt{S -> NP VP} demands \texttt{NP.agr = VP.agr}.
For a third-person subject the base form's requirement reduces to
$\textsf{num}\!:\textsf{pl}$, so the parser unifies the subject's number with
the verb's:
\[
  [\,\textsf{num}\!:\textsf{sg}\,] \ \sqcup\ [\,\textsf{num}\!:\textsf{pl}\,]
  \ =\ \top .
\]
The values conflict, unification fails, and no $S$ is built. Reading
\texttt{bark} as a noun gives \texttt{Det N N}, which no rule licenses as a
sentence, so that branch dies too. The sentence has no parse --- but it failed
on exactly one relaxable constraint, which the matching mal-rule names.

\emph{Report.}

\begin{verbatim}
  the dog bark
          ^^^^   subject-verb agreement error

  Verdict:   ungrammatical

  Violation: subject-verb agreement (number)
             subject "the dog" is 3rd person singular;
             "bark" is the non-3sg present form, which needs a
             non-3sg subject. Unifying num=sg with the required
             num=pl fails.

  Rule:      S -> NP VP   with   NP.agr = VP.agr

  Fix:       "the dog barks"   (inflect the verb for 3sg), or
             "the dogs bark"   (make the subject plural)
\end{verbatim}

This is the target behaviour: a verdict, the offending span, the violated rule,
and a concrete repair --- all falling out of a single failed unification rather
than a hand-written special case.

\section{The grammar file format}

The engine reads a grammar from a plain-text \texttt{.gram} file. The format is a
concrete syntax for the unification grammar of Section~11 --- PATR-II in spirit
(Shieber, 1986) --- with the diagnostic and mal-rule lines of Section~12 added.
There are three statement kinds: lexicon entries, phrase-structure rules, and
mal-rules. Whitespace separates tokens but is otherwise insignificant;
\texttt{\#} begins a comment running to end of line; blank lines are ignored.
The grammar of the format (\texttt{::=} defines, \texttt{|} alternates,
\texttt{\{\,\}} repeats zero or more, \texttt{[\,]} is optional, \texttt{EOL} is
end of line):

\begin{verbatim}
  grammar    ::= { lexentry | rule | malrule }

  lexentry   ::= word ":" cat { featpath "=" value } EOL

  rule       ::= production { equation }
  production ::= sym "->" sym { sym } EOL
  equation   ::= rulepath "=" ( rulepath | value ) [ diag ] EOL
  diag       ::= "!" string [ "fix:" strategy { "|" strategy } ]

  malrule    ::= "%mal" production "*ERR" string [ "*FIX" text ]

  featpath   ::= "<" feature { feature } ">"        # implicit constituent: the
                                                    # entry's own structure
  rulepath   ::= "<" sym { feature } ">"            # first sym = a constituent
                                                    # of the production; the rest
                                                    # navigate into its structure

  word       ::= token            # surface form, e.g. dog, barks
  cat, sym   ::= identifier       # category / nonterminal: N NP V Det Pron S
  feature    ::= identifier       # num pers case vagr
  value      ::= identifier       # sg pl 3 nom acc 3sg n3sg
  strategy   ::= identifier       # agree-verb nominative-pronoun ...
  string     ::= '"' { char } '"'
\end{verbatim}

The semantics the syntax does not show:
\begin{itemize}
  \item \texttt{=} is unification, not assignment (Section~5): symmetric,
        order-independent, a declared constraint rather than a store. All of a
        rule's equations are unified together when the rule applies; any
        conflict fails the application and, with it, the parse along that branch.
  \item A \texttt{rulepath} names a constituent by its category symbol, then a
        feature sequence into that constituent's structure. In a
        \texttt{lexentry} the constituent is implicit --- the entry's own
        structure --- so the path is features only.
  \item An omitted feature is unconstrained and unifies with anything; grammatical
        facts such as ``the past tense marks no agreement'' are stated by saying
        nothing (Section~11).
  \item A \texttt{!}-tagged equation carries the diagnostic reported when
        relaxing exactly that constraint rescues the parse; \texttt{fix:} names
        repair strategies. A \texttt{\%mal} rule instead recognizes a wrong
        pattern directly and always reports.
\end{itemize}

One caveat in the present design: a path identifies a constituent by category
symbol, which is unambiguous only when that symbol occurs once in the production.
A rule with a repeated category (\texttt{S -> NP NP}) would need indexed names
(\texttt{NP.1}, \texttt{NP.2}); the current fragment avoids the situation.

\begin{thebibliography}{9}
\bibitem{chomsky1956} N.\ Chomsky. \emph{Three models for the description of
language}. IRE Transactions on Information Theory, 1956.
\bibitem{shieber1985} S.\ M.\ Shieber. \emph{Evidence against the
context-freeness of natural language}. Linguistics and Philosophy, 1985.
\bibitem{joshi1985} A.\ K.\ Joshi. \emph{Tree adjoining grammars: How much
context-sensitivity is required to provide reasonable structural descriptions?}
In Natural Language Parsing, 1985.
\bibitem{vsw1994} K.\ Vijay-Shanker and D.\ J.\ Weir. \emph{The equivalence of
four extensions of context-free grammars}. Mathematical Systems Theory, 1994.
\bibitem{pullum2001} G.\ K.\ Pullum and B.\ C.\ Scholz. \emph{On the distinction
between model-theoretic and generative-enumerative syntactic frameworks}. In
LACL, 2001.
\bibitem{kasper1986} R.\ Kasper and W.\ Rounds. \emph{A logical semantics for
feature structures}. In ACL, 1986.
\bibitem{jurafsky} D.\ Jurafsky and J.\ H.\ Martin. \emph{Speech and Language
Processing}, 3rd ed.\ (draft).
\bibitem{sipser} M.\ Sipser. \emph{Introduction to the Theory of Computation}.
\bibitem{shieber1986} S.\ M.\ Shieber. \emph{An Introduction to
Unification-Based Approaches to Grammar}. CSLI, 1986.
\bibitem{pereira} F.\ C.\ N.\ Pereira and S.\ M.\ Shieber. \emph{Prolog and
Natural Language Analysis}. CSLI, 1987.
\end{thebibliography}

\end{document}
